% \n\s*\n
\subsection{Differentiation}
\tsDef{Derivative}{
    Let $f : D \to \mathbb{R}$ and assume $D$ has no isolated points.
    The \emph{difference quotient} is
    \[
        \frac{f(a+h) - f(a)}{h}.
    \]
    If the limit exists, the \emph{derivative} of $f$ at $a$ is defined by
    \[
        f'(a)
        =
        \lim_{h \to 0} \frac{f(a+h) - f(a)}{h}
        =
        \lim_{b \to a} \frac{f(b) - f(a)}{b - a}.
    \]Interpretation:
    \begin{itemize}[noitemsep]
        \item instantaneous rate of change
        \item slope of the tangent
        \item sign indicates increase/decrease
    \end{itemize}
}
\tsDef{Derivative Function}{
    Given $f : D \to \mathbb{R}$, the \emph{derivative function} is
    \(
        x \mapsto f'(x).
    \)
}
\tsDef{Differentiability}{
    A function $f$ is \emph{differentiable at $x$} if $f'(x)$ exists.If this holds for all $x$ in the domain, $f$ is called differentiable.
}
\tsLem{Differentiability Implies Continuity}{
    If $f$ is differentiable at $x_0$, then $f$ is continuous at $x_0$.
}
\tsDef{One-Sided Derivatives}{
    Derivatives
    \begin{itemize}[noitemsep]
        \item Right derivative:
              \(
                  \lim_{h \to 0^+} \frac{f(a+h) - f(a)}{h}
              \)
        \item Left derivative:
              \(
                  \lim_{h \to 0^-} \frac{f(a+h) - f(a)}{h}
              \)
    \end{itemize}
}
\tsDef{Higher Derivatives}{
    Define iteratively:
    \[
        f^{(0)} = f, \quad f^{(1)} = f', \quad f^{(2)} = f'', \quad f^{(n+1)} = (f^{(n)})'.
    \]If $f^{(n)}$ exists, $f$ is $n$-times differentiable. If all derivatives up to order $n$ are continuous, then $f \in C^n(D)$. Define
    \(
        C^\infty(D) := \bigcap_{n=0}^{\infty} C^n(D),
    \)
    called \emph{smooth functions}.
}
\tsThe{Basic Derivatives}{
    Basic Derivatives, more in the end*
    \begin{itemize}[noitemsep]
        \item Power:
              \(
                  (ax^p)' = a p x^{p-1}
              \)
        \item Exponential:
              \(
                  (a^x)' = \ln(a)\, a^x
              \)
        \item Trig.:
              \(
                  (\sin x)' = \cos x,\quad (\cos x)' = -\sin x,\quad (\tan x)' = \frac{1}{\cos^2 x}
              \)
        \item Hyperbolic:
              \(
                  (\sinh x)' = \cosh x,\quad (\cosh x)' = \sinh x
              \)
        \item Inverse:
              \(
                  (\ln x)' = \frac{1}{x},\quad (\arctan x)' = \frac{1}{1+x^2}
              \)
    \end{itemize}
}
\tsThe{Derivative Rules}{
    Derivative Rules
    \begin{itemize}[noitemsep]
        \item \textbf{Sum and Product Rule.}  
        For differentiable functions $f$ and $g$,
        \[
            (f+g)' = f' + g'
        \quad
            (fg)' = f'g + fg'.
        \]
        \item \textbf{Chain Rule.}  
        If $f$ is differentiable at $x_0$ and $g$ is differentiable at $f(x_0)$, then
        \[
            (g\circ f)'(x_0)
            =
            g'\bigl(f(x_0)\bigr)\,f'(x_0).
        \]
        \item \textbf{Quotient Rule.}  
        If $g(x_0)\neq 0$, then
        \[
            \left(\frac{f}{g}\right)'(x_0)
            =
            \frac{
                f'(x_0)g(x_0)-f(x_0)g'(x_0)
            }{
                g(x_0)^2
            }.
        \]
        \item \textbf{Derivative of the Inverse.}  
        Let $f$ be bijective and differentiable, with $f'(x_0)\neq 0$. Then
        \[
            (f^{-1})'\bigl(f(x_0)\bigr)
            =
            \frac{1}{f'(x_0)}.
        \]
    \end{itemize}
}
\tsThe{Fermat's Theorem (Necessary Condition for Extremum)}{
    If $x_0$ is a local extremum and $f$ is differentiable at $x_0$, then
    \(
        f'(x_0) = 0.
    \)
}
\tsThe{Characterization of Local Extrema}{
    If $x_0$ is a local extremum, then at least one holds:
    \begin{itemize}[noitemsep]
        \item $x_0$ is a boundary point
        \item $f$ is not differentiable at $x_0$
        \item $f'(x_0) = 0$
    \end{itemize}
}
\tsThe{Rolle's Theorem}{
    Let $f$ be continuous on $[a,b]$ and differentiable on $(a,b)$ with $f(a)=f(b)$.Then there exists $\xi \in (a,b)$ such that
    \(
        f'(\xi) = 0.
    \)
}
\tsThe{Mean Value Theorem (Lagrange)}{
    Let $f$ be continuous on $[a,b]$ and differentiable on $(a,b)$.Then there exists $\xi \in (a,b)$ such that
    \[
        f'(\xi) = \frac{f(b)-f(a)}{b-a}.
    \]
}
\tsThe{l'Hospital's Rule (0/0)}{
    If
    \(
        \lim_{x\to a} f(x)=\lim_{x\to a} g(x)=0
    \)
    and
    \(
        \lim_{x\to a} \frac{f'(x)}{g'(x)} = L,
    \)
    then
    \[
        \lim_{x\to a} \frac{f(x)}{g(x)} = L.
    \]
}
\tsThe{l'Hospital's Rule ($\infty/\infty$)}{
    If
    \(
        \lim_{x\to a} |f(x)|=\lim_{x\to a} |g(x)|=\infty
    \)
    and
    \(
        \lim_{x\to a} \frac{f'(x)}{g'(x)} = L,
    \)
    then
    \[
        \lim_{x\to a} \frac{f(x)}{g(x)} = L.
    \]
}
\tsThe{l'Hospital at Infinity}{
    If
    \(
        \lim_{x\to\infty} f(x), g(x) \in \{0,\infty\}
    \)
    and
    \(
        \lim_{x\to\infty} \frac{f'(x)}{g'(x)} = L,
    \)
    then
    \[
        \lim_{x\to\infty} \frac{f(x)}{g(x)} = L.
    \]
}
\tsThe{Monotonicity Criterion (First Derivative Test)}{
    \(
        f'(x) \ge 0 \iff f \text{ is increasing}
    \)
}
\tsCor{Constant Functions}{
    $f$ is constant iff
    \(
        f'(x) = 0 \quad \forall x.
    \)
}
\tsThe{Convexity via First Derivative}{
    $f$ is convex iff $f'$ is increasing.
}
\tsCor{Convexity via Second Derivative}{
    If $f$ is twice differentiable:
    \begin{itemize}[noitemsep]
        \item \(f''(x) \ge 0 \iff f \text{ is convex}.\)
        \item \(f''(x) \leq 0 \Longleftrightarrow f \text{ is concave}\)
    \end{itemize}
}
\tsThe{Second Derivative Test}{
Let $f$ be twice differentiable near $x_0$ and suppose
\(
f'(x_0)=0.
\)
Then:
\(
f''(x_0)>0
\quad \Rightarrow \quad
x_0 \text{ is a local minimum},
\)
and
\(
f''(x_0)<0
\quad \Rightarrow \quad
x_0 \text{ is a local maximum}.
\)
}
\subsection{Taylor Approximation and Power Series}
\tsIdea{Local Approximation by Polynomials}{
In a sufficiently small neighborhood of a point $x_0$, the graph of a differentiable function $f$ behaves almost like its tangent at $x_0$. Hence, the function can locally be approximated by its tangent.
}
\tsDef{Linear Approximation / Tangent Approximation}{
If $f$ is differentiable at $x_0$, then
\(
f(x)
\approx
f(x_0) + f'(x_0)(x-x_0).
\)
The expression
\(
P_1(x)
=
f(x_0) + f'(x_0)(x-x_0)
\)
is called the \emph{linearization}, \emph{tangent approximation}, or \emph{first Taylor polynomial} of $f$ at $x_0$.
}
\tsThe{Taylor's Theorem (with Lagrange Remainder)}{
Let $f$ be infinitely differentiable on an open interval $I$ containing $x_0$.
Then for every $x \in I$,
\[
f(x)
=
f(x_0)
+
f'(x_0)(x-x_0)
+\cdots+
\frac{1}{n!}f^{(n)}(x_0)(x-x_0)^n,
\]
for some $c$ between $x_0$ and $x$.
Equivalently,
\[
f(x)
=
\sum_{k=0}^{n}
\frac{f^{(k)}(x_0)}{k!}(x-x_0)^k
+
R_n(x),
\]
where
\(
R_n(x)
=
\frac{1}{(n+1)!}
f^{(n+1)}(c)(x-x_0)^{n+1}.
\)
The polynomial
\[
P_n(x)
=
\sum_{k=0}^{n}
\frac{f^{(k)}(x_0)}{k!}(x-x_0)^k
\]
is called the \emph{$n$-th Taylor polynomial}.
}
\tsIdea{Taylor Approximation and Power Series}{
Taylor polynomials provide increasingly accurate approximations of many functions.
For many functions, taking infinitely many terms recovers the original function.
}
\tsDef{Taylor Series}{
For a function $f$, the series
\(
\sum_{k=0}^{\infty}
\frac{f^{(k)}(a)}{k!}(x-a)^k
\)
is called the \emph{Taylor series} (or Taylor expansion) of $f$ around the point $a$.
}
\tsThe{Important Taylor Series}{
The following series converge for all $x$:
\[
\sin(x)
=
\sum_{n=0}^{\infty}
(-1)^n
\frac{x^{2n+1}}{(2n+1)!}
=
x - \frac{x^3}{3!} + \frac{x^5}{5!} - \cdots
\]
\[
\cos(x)
=
\sum_{n=0}^{\infty}
(-1)^n
\frac{x^{2n}}{(2n)!}
=
1 - \frac{x^2}{2!} + \frac{x^4}{4!} - \cdots
\]
\[
e^x
=
\sum_{n=0}^{\infty}
\frac{x^n}{n!}
=
1 + x + \frac{x^2}{2!} + \frac{x^3}{3!} + \cdots
\]
}
\tsThe{Further Important Power Series (Geometric, Binomial)}{
The following series converge for
\(
-1 < x < 1.
\)
\[
\ln(1+x)
=
\sum_{n=1}^{\infty}
(-1)^{n-1}
\frac{x^n}{n}
=
x - \frac{x^2}{2} + \frac{x^3}{3} - \frac{x^4}{4} + \cdots
\]
\[
\frac{1}{1-x}
=
\sum_{n=0}^{\infty} x^n
=
1 + x + x^2 + x^3 + \cdots
\]
\[
(1+x)^p
=
\sum_{n=0}^{\infty}
\binom{p}{n}x^n
=
1 + px + \frac{p(p-1)}{2!}x^2 + \cdots
\]
}
\tsIdea{Operations on Taylor Series}{
New Taylor series can be constructed from known ones using substitution, multiplication, differentiation, integration.
}
\tsThe{Cauchy Product of Power Series}{
Let
\[
f(x)=\sum_{n=0}^{\infty} a_nx^n,
\qquad
g(x)=\sum_{n=0}^{\infty} b_nx^n,
\]
with $x$ inside both convergence intervals.
Then
\[
f(x)g(x)
=
\sum_{n=0}^{\infty}
\left(
\sum_{k=0}^{n} a_k b_{n-k}
\right)x^n.
\]
}
\tsThe{Termwise Differentiation}{
Let
\(
f(x)=\sum_{n=0}^{\infty} a_nx^n
\)
inside its convergence interval. Then $f$ is differentiable and
\[
f'(x)
=
\sum_{n=1}^{\infty}
n a_n x^{n-1}.
\]
}
\tsThe{Termwise Integration}{
Let
\(
f(x)=\sum_{n=0}^{\infty} a_nx^n
\)
inside its convergence interval $I$. Then for every interval $[a,b]\subseteq I$,
\[
\int_a^b f(x)\,dx
=
\sum_{n=0}^{\infty}
\int_a^b a_nx^n\,dx.
\]
}